% \section{Related Work}
% \label{sec:related}

% We summarize some related work in the literature
% , which can be classified into two categories, i.e., the CDP methods and the
% LDP methods.

% \subsection{CDP Methods}
% \paragraph{Event-level DP}
% Dwork et al.~ first introduced event-level DP for
% continuously releasing statistics like counts and histograms~\cite{Dwork2006}.
% Hierarchical tree structures are widely used for range queries. In such trees, leaf nodes store the noisy counts for each timestamp, while internal nodes store noisy sums over the intervals they cover.
% For fixed-length binary streams, a binary tree structure was used to aggregate
% counts~\cite{Dwork2010}, and later extended to infinite streams by Chan et
% al.~through consistency constraints~\cite{Chan2011}.
% To further reduce noise in sparse regions, Dwork et al.~proposed adaptive partitioning
% based on thresholds~\cite{Dwork2015}.
% However, this method is limited to post-partition release and lacks real-time
% applicability.
% Chen et al.~\cite{Chen2017} improved flexibility by modifying monitored events.
% Cao et al.~\cite{Cao2015} introduced group-based histogram publishing, adding
% Laplace noise to similar time slots.
% Although these techniques support numerical data, they require raw data access
% and thus cannot be directly applied under LDP.
% Bao et al.~\cite{Bolot2013} addressed sliding-window predicate sum queries using
% decay models.
% \paragraph{User-level DP}
% User-level DP frameworks focus on continuously publishing aggregate statistics.
% Fan et al.~\cite{Fan2013a} proposed FAST, a real-time
% aggregation framework using adaptive sampling and filtering.
% It was extended to 2D monitoring via spatial partitioning with
% quadtrees~\cite{Fan2013}.
% Rastogi et al.~\cite{Rastogi2010} developed a method based on discrete Fourier
% transform (DFT), though it is suited only for offline analysis. The analysis in \cite{2023userdp} reveals that on infinite streams, achieving user-level DP forces error to grow without bound as the maximum per-user prefix contribution increases over time. While they control privacy by truncating per-user contributions, this approach in turn introduces bias.
% \paragraph{$w$-event-level DP} 
% Kellaris et al.~\cite{Kellaris2014} introduced the $w$-event DP model to protect
% sequences of events over sliding windows and proposed BA/BD mechanisms.
% However, they rely on raw data to compute similarity, exposing them to inference
% attacks.
% Cao et al.~\cite{Cao2015} extended BD to support variable-length trajectories,
% using a greedy strategy to match current data with historical outputs.
% This resulted in uneven privacy budget distribution.
% Du et al.~\cite{du2025} extended BD/BA to the personalized $w$-event privacy setting, where different users have heterogeneous privacy requirements, and proposed PBD and PBA.
% Wang et al.~\cite{Wang2019} proposed E-RescueDP, which adaptively allocates
% privacy budgets using an RNN to support real-time release.
% Li et al.~\cite{Li2025} proposed SPAS, which predicts future data variation to
% adaptively determine data sampling and privacy budget allocation under $w$-event
% DP.

% \subsection{LDP Methods}
% Most existing LDP works focus on single-value counting and frequency estimation
% of static data, with only a few studies focusing on stream data analysis.
% Memoization-based techniques~\cite{Ding2017, Erlingsson2014, Arcolezi2021} were
% proposed to offer longitudinal LDP guarantees.
% Joseph et al.~\cite{Joseph2018} estimated stream means by having users compare
% their local averages against the last released value and vote on updates.
% Although satisfying event-level LDP, this method assumes Bernoulli input and
% temporal independence, limiting its generality.
% THRESH further assumes that the number of global updates is bounded by
% distribution changes, making it unsuitable for infinite streams.
% Wang et al.~\cite{Wang2021} proposed an event-level LDP framework for interval
% sum estimation using a hybrid mechanism with thresholding.
% However, it directly extends CDP-style hierarchies, offering limited protection
% for unbounded streams.
% Bao et al.~\cite{Bao2021} leveraged autocorrelation to reduce noise via an
% analytic Gaussian mechanism, but their method applies only to finite data and
% achieves approximate $(\epsilon, \delta)$-LDP under periodic budget renewal.
% Erlingsson et al.~\cite{Erlingsson2019} introduced a shuffling model for
% correlated time series under user-level LDP, but their approach assumes
% integer-valued inputs with bounded updates, and the hierarchical design
% restricts scalability to infinite streams.


\section{Related Work}
\label{sec:related}

We summarize some related work in the literature.
% , which can be classified into two categories, i.e., the CDP methods and the
% LDP methods.

\paragraph{CDP Methods}
In the centralized setting, Dwork et al.~introduced event-level DP for
continuously releasing statistics like counts and histograms~\cite{Dwork2006}.
Hierarchical tree structures are widely used for temporal range queries under CDP. In such trees, leaf nodes store the noisy counts for each timestamp, while internal nodes store noisy sums over the intervals they cover.
For fixed-length binary streams, a binary tree structure was used to aggregate
counts~\cite{Dwork2010}, and later extended to infinite streams by Chan et
al.~through consistency constraints~\cite{Chan2011}.
To further reduce noise in sparse regions, Dwork proposed adaptive partitioning
based on thresholds~\cite{Dwork2015}.
However, this method is limited to post-partition release and lacks real-time
applicability.
Chen et al.~\cite{Chen2017} improved flexibility by modifying monitored events.
Cao et al.~\cite{Cao2015} introduced group-based histogram publishing, adding
Laplace noise to similar time slots.
Although these techniques support numerical data, they require raw data access
and thus cannot be directly applied under LDP.
Bao et al.~\cite{Bolot2013} addressed sliding-window predicate sum queries using
decay models.

For user-level DP, Fan et al.~\cite{Fan2013a} proposed FAST, a real-time
aggregation framework using adaptive sampling and filtering.
It was extended to 2D monitoring via spatial partitioning with
quadtrees~\cite{Fan2013}.
Rastogi et al.~\cite{Rastogi2010} developed a method based on discrete Fourier
transform (DFT), though it is suited only for offline analysis. {The analysis in \cite{2023userdp} reveals that on infinite streams, achieving user-level DP forces error to grow without bound as the maximum per-user prefix contribution increases over time. While they control privacy by truncating per-user contributions, this approach in turn introduces bias.

Kellaris et al.~\cite{Kellaris2014} introduced the $w$-event DP model to protect
sequences of events over sliding windows and proposed BA/BD mechanisms.
However, they rely on raw data to compute similarity, exposing them to inference
attacks.
Cao et al.~\cite{Cao2015} extended BD to support variable-length trajectories,
using a greedy strategy to match current data with historical outputs.
This resulted in uneven privacy budget distribution.
Du et al.~\cite{du2025} extended BD/BA to the personalized $w$-event privacy setting, where different users have heterogeneous privacy requirements, and proposed PBD and PBA.
Wang et al.~\cite{Wang2019} proposed E-RescueDP, which adaptively allocates
privacy budgets using an RNN to support real-time release.
Li et al.~\cite{Li2025} proposed SPAS, which predicts future data variation to
adaptively determine data sampling and privacy budget allocation under $w$-event
DP.

\paragraph{LDP Methods}
Most existing LDP works focus on single-value counting and frequency estimation
of static data, with only a few studies focusing on stream data analysis.
Memoization-based techniques~\cite{Ding2017, Erlingsson2014, Arcolezi2021} were
proposed to offer longitudinal LDP guarantees.
Joseph et al.~\cite{Joseph2018} estimated stream means by having users compare
their local averages against the last released value and vote on updates.
Although satisfying event-level LDP, this method assumes Bernoulli input and
temporal independence, limiting its generality.
THRESH further assumes that the number of global updates is bounded by
distribution changes, making it unsuitable for infinite streams.
Wang et al.~\cite{Wang2021} proposed an event-level LDP framework for interval
sum estimation using a hybrid mechanism with thresholding.
However, it directly extends CDP-style hierarchies, offering limited protection
for unbounded streams.
Bao et al.~\cite{Bao2021} leveraged autocorrelation to reduce noise via an
analytic Gaussian mechanism, but their method applies only to finite data and
achieves approximate $(\epsilon, \delta)$-LDP under periodic budget renewal.
Erlingsson et al.~\cite{Erlingsson2019} introduced a shuffling model for
correlated time series under user-level LDP, but their approach assumes
integer-valued inputs with bounded updates, and the hierarchical design
restricts scalability to infinite streams.
